\scanpage{138}
\mspage{137}
\margin{- 16 -}
\margin{n° 505}
\uline{P. 129}, ligne 3 du haut, remplacer ``métrisable'' par ``séparé''. Ligne 19 du bas, remplacer ``\uline{métrisable}'' par ``\uline{séparé}''.

Lignes 2 à 10 du bas, remplacer par:

Considérons en effet l'application surjective $f:\{0,1\}^{\uuline{N}}\to I$ qui, à tout point $(\unsure{\epsilon|e}_n)_{n\ge0}$ avec $\unsure{\epsilon|e}_n=0$ ou $\unsure{\epsilon|e}_n=1$, fait correspondre le nombre $\sum_{n=0}^{\infty}\unsure{\epsilon|e}_n2^{-n}$ ayant $(\unsure{\epsilon|e}_n)$ pour développement dyadique. Soit $D$ l'ensemble dénombrable dans $\{0,1\}^{\uuline{N}}$ formé des suites autres que la suite $0$, n'ayant qu'un nombre fini de termes $\ne0$; $P=\{0,1\}^{\uuline{N}}-D$ est un espace polonais éparpillé (n° 1, th. 1), et la restriction de $f$ à $P$ est une \uline{bijection} continue de $\unsure{P'|P}$ sur $I$ (chap. IV, § 8, n° 6).

\uline{P. 130}, la déf. 7 devient déf. 8; après la ligne 4 du haut, insérer:

COROLLAIRE.-- \uline{Toute somme dénombrable d'espaces lusiniens est un espace lusinien.}

\uline{Toute limite projective dénombrable} (chap. I, § 4, n° 4) \uline{d'espaces lusiniens est un espace lusinien.}

La première assertion découle aussitôt de la déf. 5 et du n° 1, prop. 1; la seconde résulte de ce qu'une limite projective dénombrable d'espaces séparés est un sous-espace fermé du produit (dénombrable) de ces espaces, et de ce qu'on a vu au cours de la démonstration de la prop. 11.

${}^*$On peut montrer que les espaces de distributions $D'(\Omega)$, $E'(\Omega)$, etc. sur un ouvert $\Omega$ de $\uuline{R}^{n}$, munis de la topologie forte, sont des espaces lusiniens.${}_*$

Lignes 16 à 20 du haut, remplacer par: ensemble de la topologie \uline{limite projective} des topologies discrètes sur les $C_n$ (chap. I, § 4, n° 4).

Ligne 10 du bas, remplacer ``\uline{fermées}'' par ``\uline{ouvertes}''. Ligne 6 du bas, ajouter à la fin: et contienne les adhérences $\overline{\varphi_{n+1}(c')}$ des $\varphi_{n+1}(c')$ dans $E$ (pour
